\subsection{长度与约化分解}

定义 1. 设 \(w \in W\)。\(w\) 的长度（相对于 \(S\)），记为 \(l_{S}(w)\) 或简称为 \(l(w)\)，是满足 \(q \geqslant 0\) 且 \(w\) 是由 \(q\) 个 \(S\) 中元素组成的序列的乘积的最小整数。\(w\) 的一个约化分解（相对于 \(S\)）是任意序列 \(s = (s_1, \ldots, s_q)\)，其元素属于 \(S\)，满足 \(w = s_1 \ldots s_q\) 且 \(q = l(w)\)。

因此，1 是长度为 0 的唯一元素，而 \(S\) 由长度为 1 的元素组成。

命题 1. 设 w 和 w' 属于 W。我们有以下公式：

\[
l (w w ^ {\prime}) \leqslant l (w) + l (w ^ {\prime}),\tag{1}
\]

\[
l (w ^ {- 1}) = l (w),\tag{2}
\]

\[
\left| l (w) - l \left(w ^ {\prime}\right) \right| \leqslant l \left(w w ^ {\prime - 1}\right).\tag{3}
\]

设 \((s_{1},\ldots,s_{p})\) 和 \((s_{1}^{\prime},\ldots,s_{q}^{\prime})\) 分别是 w 和 \(w^{\prime}\) 的约化分解。因此 \(l(w)=p\) 且 \(l(w^{\prime})=q\) 。由于 \(ww^{\prime}=s_{1}\ldots s_{p}s_{1}^{\prime}\ldots s_{q}^{\prime}\)，有 \(l(ww^{\prime})\leqslant p+q\)，这就证明了 (1)。由于 \(S=S^{-1}\) 且 \(w^{-1}=s_{p}^{-1}\ldots s_{1}^{-1}\)，有 \(l(w^{-1})\leqslant p=l(w)\)。用 \(w^{-1}\) 代替 w，得到相反的不等式 \(l(w)\leqslant l(w^{-1})\)，从而证明 (2)。在 (1) 和 (2) 中以 \(ww^{\prime-1}\) 代替 w，得到关系

\[
l (w) - l \left(w ^ {\prime}\right) \leqslant l \left(w w ^ {\prime - 1}\right),\tag{4}
\]

\[
l \left(w w ^ {\prime - 1}\right) = l \left(w ^ {\prime} w ^ {- 1}\right).\tag{5}
\]

在 (4) 中交换 \(w\) 和 \(w'\)，由 (5) 得到 \(l(w') - l(w) \leqslant l(ww'^{-1})\)，从而证明 (3)。

推论。设 \(\mathbf{s} = (s_{1}, \ldots, s_{p})\) 和 \(\mathbf{s}' = (s_{1}', \ldots, s_{q}')\) 是 S 中元素的两个序列，满足 \(w = s_{1} \ldots s_{p}\) 且 \(w' = s_{1}' \ldots s_{q}'\)。如果序列 \((s_{1}, \ldots, s_{p}, s_{1}', \ldots, s_{q}')\) 是 \(ww'\) 的一个约化分解，则 s 是 w 的一个约化分解，而 \(s'\) 是 \(w'\) 的一个约化分解。

根据假设，\(l(w) \leqslant p\)，\(l(w') \leqslant q\)，且 \(l(ww') = p + q\)。由 (1)，必有 \(l(w) = p\) 且 \(l(w') = q\)，于是得到该推论。

备注。由公式 (1) 和 (2)，公式 \(d(w, w') = l(ww'^{-1})\) 在 W 上定义了一个距离，该距离在右平移下不变。

